\section{把作业当成课程的延伸}

这门课的作业不是脱离讲义的单独项目，而是把 lecture 中的抽象变成手上能跑的系统。
如果把五个作业放在一起看，它们其实构成了一条从“最小语言模型”走到“可用模型”的路线。

\begin{table}[H]
\centering
\begin{tabular}{p{2.1cm}p{3.1cm}p{4.3cm}p{3.5cm}}
\toprule
\textbf{作业} & \textbf{对应模块} & \textbf{你会真正做什么} & \textbf{最常见失败点} \\
\midrule
HW1 & Basics & tokenizer, Transformer, training loop & 实现正确但太慢 \\
HW2 & Systems & Triton kernel, DDP, state sharding & 只对 correctness，不对性能负责 \\
HW3 & Scaling laws & 拟合和外推 scaling 曲线 & 小实验不稳，外推不可信 \\
HW4 & Data & 采集、转换、过滤、去重 & 数据质量差，指标看不出来 \\
HW5 & Alignment & SFT, DPO, GRPO & 只看 loss，不看偏好和安全目标 \\
\bottomrule
\end{tabular}
\caption{作业不是“额外负担”，而是课程主线的工程化版本}
\end{table}

\begin{warningbox}{做作业时最危险的心态}
只把作业看成“把分数拿到手”，而不去看它在系统层面对应什么设计问题。这样做完以后，知识不会真正积累。
\end{warningbox}

\begin{knowledgebox}{做完这门课后你应该能回答的问题}
\begin{itemize}
    \item 为什么这个 tokenizer 不够好？
    \item 为什么这个模型在小规模上有效，在大规模上未必最优？
    \item 为什么系统优化能改变训练策略，而不仅是加速实现？
    \item 为什么数据和对齐会决定模型最终是否“有用”？
\end{itemize}
\end{knowledgebox}

